% !TEX root = ../main.tex

% ==================================================
\section{Discussion and Limitations}
\label{sec:discussion}
% ==================================================

\VLMemCodec{} is designed around efficient visual transmission in Edge--Cloud VSLAM rather than generic video delivery. Its VLMem payload keeps visual anchors and language cues as compact Vision--Language Memory, and the reconstructed sequence is evaluated primarily by downstream VSLAM trajectory quality. This task-driven framing should not be interpreted as a closed-form optimization of the full VSLAM pipeline; instead, it provides a practical coding design whose trajectory quality is assessed through the cloud-side VSLAM estimator.

The results also show that image-fidelity metrics and downstream VSLAM trajectory quality are not fully aligned. LPIPS, SSIM, and FID are useful auxiliary metrics because they characterize complementary aspects of reconstructed image appearance, but they do not necessarily move in the same direction as APE RMSE. In the MGDR ablation, decoder refinement implemented through LoRA-based low-rank increments reduces REC APE in the reported setting, while the auxiliary image-fidelity metrics do not improve uniformly.

The current prototype demonstrates the communication benefit of the compact VLMem payload, but cloud-side world-model reconstruction remains the major runtime bottleneck in the present implementation. This separates the communication benefit from the current reconstruction-runtime bottleneck. The compact transmitted representation reduces payload size, while final end-to-end latency also depends on the world-model decoder, accelerator memory, batching and parallel serving, and inference optimization. The present system should therefore be viewed as a prototype for validating the efficient visual transmission formulation in Edge--Cloud VSLAM, not as a fully real-time optimized cloud service.

The cloud-side world-model decoder also introduces a reconstruction limitation. Although visual anchors and language cues constrain the reconstruction, the reconstructed frames may still contain hallucinations, local geometric inconsistency, viewpoint drift, dynamic-object artifacts, or texture changes. Therefore, reconstructed frames should be treated as task-usable reconstructions for downstream VSLAM evaluation, not physically accurate observations.

The edge-side encoder uses lightweight motion and semantic event cues rather than precise dense motion estimation. Phase correlation provides coarse language-level motion conditions, not full optical flow. Future work will focus on faster world-model serving, model compression or distillation, and parallel reconstruction of VLMem units; stronger geometry-aware constraints, closed-loop SLAM feedback, and uncertainty-aware decoding; and task-specific lightweight decoders for efficient visual transmission in Edge--Cloud VSLAM.

% ==================================================
\section{Conclusion}
\label{sec:conclusion}
% ==================================================

This paper presented \VLMemCodec, a Vision--Language Memory coding framework with world-model reconstruction for efficient visual transmission in Edge--Cloud VSLAM. The proposed framework represents dense image streams as a VLMem payload with visual anchors and language cues, and reconstructs task-usable visual sequences on the cloud using a vision--language conditioned world-model decoder refined through MGDR.

The experimental results show a favorable compression--trajectory tradeoff: the transmitted VLMem payload substantially reduces visual data, REC APE remains close to ORI APE on average, phase correlation provides low-cost motion cues for the edge encoder, and MGDR reduces REC APE in the decoder ablation. The results also indicate that auxiliary image-fidelity metrics are not fully aligned with VSLAM trajectory quality. Future work will further improve geometric consistency and closed-loop VSLAM-aware reconstruction.
